\documentclass[11pt,a4paper]{article}

%=========================================================
% PACKAGES
%=========================================================
\usepackage[utf8]{inputenc}
\usepackage[T1]{fontenc}
\usepackage[french]{babel}
\usepackage[a4paper,margin=2.15cm]{geometry}
\usepackage{lmodern}
\usepackage{amsmath,amssymb,amsthm,mathtools}
\usepackage{fancyhdr}
\usepackage{lastpage}
\usepackage{enumitem}
\usepackage{array,tabularx}
\usepackage{tikz}
\usepackage[most]{tcolorbox}
\usepackage{hyperref}
\usepackage{xcolor}

%=========================================================
% MISE EN PAGE
%=========================================================
\setlength{\headheight}{14pt}
\pagestyle{fancy}
\fancyhf{}
\fancyhead[L]{Mathématiques CPGE}
\fancyhead[C]{Diagonalisation et trigonalisation}
\fancyhead[R]{Page \thepage/\pageref{LastPage}}
\fancyfoot[C]{Cours complet -- Réduction des endomorphismes}

\hypersetup{
    colorlinks=true,
    linkcolor=blue!50!black,
    urlcolor=blue!50!black
}

\setlist[itemize]{label=\(\triangleright\)}
\setlist[enumerate]{label=\textbf{\arabic*.}}

%=========================================================
% COULEURS ET BOÎTES
%=========================================================
\definecolor{bleufonce}{RGB}{25,65,130}
\definecolor{bleuclair}{RGB}{235,243,255}
\definecolor{vertfonce}{RGB}{30,120,70}
\definecolor{vertclair}{RGB}{235,250,240}
\definecolor{rougefonce}{RGB}{170,40,40}
\definecolor{rougeclair}{RGB}{255,238,238}
\definecolor{orangefonce}{RGB}{190,110,20}
\definecolor{orangeclair}{RGB}{255,246,230}
\definecolor{violetfonce}{RGB}{95,55,145}
\definecolor{violetclair}{RGB}{246,241,255}

\tcbset{
    enhanced,
    sharp corners,
    boxrule=0.75pt,
    before skip=8pt,
    after skip=8pt
}

\newtcolorbox{definitionbox}[1][]{
    colback=bleuclair,
    colframe=bleufonce,
    title=\textbf{Définition #1}
}

\newtcolorbox{theorembox}[1][]{
    colback=vertclair,
    colframe=vertfonce,
    title=\textbf{Théorème / Propriété #1}
}

\newtcolorbox{methodbox}[1][]{
    colback=orangeclair,
    colframe=orangefonce,
    title=\textbf{Méthode #1}
}

\newtcolorbox{retenirbox}[1][]{
    colback=violetclair,
    colframe=violetfonce,
    title=\textbf{À retenir #1}
}

\newtcolorbox{erreurbox}[1][]{
    colback=rougeclair,
    colframe=rougefonce,
    title=\textbf{Erreur classique #1}
}

\newtcolorbox{exercicebox}[1][]{
    colback=white,
    colframe=black!60,
    title=\textbf{Exercice #1}
}

\newtcolorbox{correctionbox}[1][]{
    colback=gray!5,
    colframe=black!65,
    title=\textbf{Correction #1}
}

%=========================================================
% COMMANDES
%=========================================================
\newcommand{\K}{\mathbb K}
\newcommand{\R}{\mathbb R}
\newcommand{\C}{\mathbb C}
\newcommand{\N}{\mathbb N}
\newcommand{\M}{\mathcal M}
\newcommand{\Lcal}{\mathcal L}
\newcommand{\Vect}{\operatorname{Vect}}
\newcommand{\Ker}{\operatorname{Ker}}
\newcommand{\Imf}{\operatorname{Im}}
\newcommand{\rg}{\operatorname{rg}}
\newcommand{\tr}{\operatorname{tr}}
\newcommand{\Id}{\operatorname{Id}}
\newcommand{\Sp}{\operatorname{Sp}}
\newcommand{\diag}{\operatorname{diag}}
\newcommand{\Mat}{\operatorname{Mat}}

\newtheorem{prop}{Proposition}[section]
\newtheorem{lemme}[prop]{Lemme}
\newtheorem{corollaire}[prop]{Corollaire}

%=========================================================
% DOCUMENT
%=========================================================
\begin{document}

\begin{center}
    {\Huge \textbf{Diagonalisation et trigonalisation}}\\[2mm]
    {\Large Réduire un endomorphisme pour comprendre son action}\\[2mm]
    {\large Cours complet pour classes préparatoires scientifiques}\\[2mm]
    \rule{0.88\linewidth}{0.5pt}
\end{center}

\vspace{0.4cm}

\tableofcontents

\newpage

%=========================================================
\section*{Introduction générale}
\addcontentsline{toc}{section}{Introduction générale}

Dans les chapitres précédents, on a introduit les matrices, le calcul matriciel, les polynômes, les valeurs propres et les sous-espaces propres.  
La question centrale de ce chapitre est la suivante :

\begin{center}
\textit{Peut-on choisir une base dans laquelle l'endomorphisme prend une forme simple ?}
\end{center}

Soit \(E\) un \(\K\)-espace vectoriel de dimension finie \(n\), et soit :
\[
u\in\Lcal(E).
\]
Si l'on choisit une base de \(E\), l'endomorphisme \(u\) est représenté par une matrice.  
Changer de base revient à remplacer cette matrice par une matrice semblable :
\[
A' = P^{-1}AP.
\]
La réduction consiste à chercher une base dans laquelle la matrice est plus simple.

Deux formes sont particulièrement importantes :

\begin{itemize}
    \item une forme diagonale :
    \[
    D=
    \begin{pmatrix}
    \lambda_1&0&\cdots&0\\
    0&\lambda_2&\cdots&0\\
    \vdots&\vdots&\ddots&\vdots\\
    0&0&\cdots&\lambda_n
    \end{pmatrix};
    \]
    \item une forme triangulaire supérieure :
    \[
    T=
    \begin{pmatrix}
    \lambda_1&*&\cdots&*\\
    0&\lambda_2&\cdots&*\\
    \vdots&\vdots&\ddots&\vdots\\
    0&0&\cdots&\lambda_n
    \end{pmatrix}.
    \]
\end{itemize}

\begin{retenirbox}
Diagonaliser, c'est trouver une base de vecteurs propres.  
Trigonaliser, c'est trouver une base dans laquelle l'endomorphisme est représenté par une matrice triangulaire supérieure.
\end{retenirbox}

La diagonalisation est la situation idéale.  
La trigonalisation est une réduction plus faible, mais souvent possible même lorsque la diagonalisation échoue.

%=========================================================
\section{Pré-requis essentiels}

\subsection{Matrices semblables}

\begin{definitionbox}[Matrices semblables]
Deux matrices \(A,B\in\M_n(\K)\) sont dites semblables s'il existe une matrice inversible \(P\in GL_n(\K)\) telle que :
\[
B=P^{-1}AP.
\]
\end{definitionbox}

\begin{retenirbox}
Deux matrices semblables représentent le même endomorphisme dans deux bases différentes.
\end{retenirbox}

\begin{theorembox}[Invariants par similitude]
Si \(A\) et \(B\) sont semblables, alors :
\[
\chi_A=\chi_B,
\qquad
\Sp(A)=\Sp(B),
\qquad
\det(A)=\det(B),
\qquad
\tr(A)=\tr(B).
\]
\end{theorembox}

\begin{proof}
Supposons :
\[
B=P^{-1}AP.
\]
Alors :
\[
XI_n-B=P^{-1}(XI_n-A)P.
\]
Donc :
\[
\chi_B(X)=\det(XI_n-B)
=
\det(P^{-1})\det(XI_n-A)\det(P)
=
\chi_A(X).
\]
Les valeurs propres étant les racines du polynôme caractéristique, les spectres sont égaux.

Pour le déterminant :
\[
\det(B)=\det(P^{-1})\det(A)\det(P)=\det(A).
\]
Pour la trace, on utilise la propriété :
\[
\tr(P^{-1}AP)=\tr(APP^{-1})=\tr(A).
\]
\end{proof}

\subsection{Valeurs propres et sous-espaces propres}

\begin{definitionbox}[Valeur propre]
Un scalaire \(\lambda\in\K\) est une valeur propre de \(u\in\Lcal(E)\) s'il existe \(x\neq0\) tel que :
\[
u(x)=\lambda x.
\]
\end{definitionbox}

\begin{definitionbox}[Sous-espace propre]
Le sous-espace propre associé à \(\lambda\) est :
\[
E_\lambda(u)=\Ker(u-\lambda\Id_E).
\]
\end{definitionbox}

\begin{theorembox}
On a :
\[
\lambda\in\Sp(u)
\Longleftrightarrow
E_\lambda(u)\neq\{0\}.
\]
\end{theorembox}

\begin{proof}
Par définition, \(\lambda\) est valeur propre si et seulement s'il existe \(x\neq0\) tel que :
\[
u(x)=\lambda x.
\]
Cela équivaut à :
\[
(u-\lambda\Id_E)(x)=0,
\]
avec \(x\neq0\), donc à :
\[
E_\lambda(u)\neq\{0\}.
\]
\end{proof}

\subsection{Polynôme caractéristique}

\begin{definitionbox}[Polynôme caractéristique]
Pour \(A\in\M_n(\K)\), on définit :
\[
\chi_A(X)=\det(XI_n-A).
\]
Pour un endomorphisme \(u\), on définit \(\chi_u\) comme le polynôme caractéristique d'une matrice de \(u\) dans une base quelconque.
\end{definitionbox}

\begin{theorembox}
Les valeurs propres de \(A\) sont exactement les racines de \(\chi_A\) dans \(\K\).
\[
\lambda\in\Sp(A)
\Longleftrightarrow
\chi_A(\lambda)=0.
\]
\end{theorembox}

\begin{proof}
On a :
\[
\lambda\in\Sp(A)
\Longleftrightarrow
A-\lambda I_n \text{ n'est pas inversible}.
\]
Cela équivaut à :
\[
\det(A-\lambda I_n)=0.
\]
Or :
\[
\chi_A(\lambda)=\det(\lambda I_n-A)=(-1)^n\det(A-\lambda I_n).
\]
Donc :
\[
\chi_A(\lambda)=0
\Longleftrightarrow
\det(A-\lambda I_n)=0.
\]
\end{proof}

%=========================================================
\section{Diagonalisation : définition et interprétation}

\subsection{Endomorphisme diagonalisable}

\begin{definitionbox}[Endomorphisme diagonalisable]
Un endomorphisme \(u\in\Lcal(E)\) est dit diagonalisable s'il existe une base de \(E\) formée de vecteurs propres de \(u\).
\end{definitionbox}

\subsection{Matrice diagonalisable}

\begin{definitionbox}[Matrice diagonalisable]
Une matrice \(A\in\M_n(\K)\) est dite diagonalisable s'il existe une matrice inversible \(P\in GL_n(\K)\) et une matrice diagonale \(D\in\M_n(\K)\) telles que :
\[
A=PDP^{-1}.
\]
\end{definitionbox}

\begin{retenirbox}
Dire que \(A\) est diagonalisable signifie que \(A\) est semblable à une matrice diagonale.
\end{retenirbox}

\subsection{Lien entre les deux points de vue}

\begin{theorembox}
Soit \(u\in\Lcal(E)\), et soit \(A\) la matrice de \(u\) dans une base quelconque.  
Alors :
\[
u\text{ est diagonalisable}
\Longleftrightarrow
A\text{ est diagonalisable}.
\]
\end{theorembox}

\begin{proof}
Supposons \(u\) diagonalisable.  
Il existe une base :
\[
\mathcal B'=(v_1,\dots,v_n)
\]
formée de vecteurs propres de \(u\).  
Alors :
\[
u(v_i)=\lambda_i v_i.
\]
Dans cette base, la matrice de \(u\) est :
\[
D=\diag(\lambda_1,\dots,\lambda_n).
\]
Si \(P\) est la matrice de passage de \(\mathcal B'\) vers la base initiale, on a :
\[
A=PDP^{-1}.
\]
Donc \(A\) est diagonalisable.

Réciproquement, si :
\[
A=PDP^{-1},
\]
alors les colonnes de \(P\) forment une base de \(\K^n\), et l'égalité :
\[
AP=PD
\]
montre que ces colonnes sont des vecteurs propres de \(A\).  
Donc l'endomorphisme associé est diagonalisable.
\end{proof}

\begin{methodbox}[Lire une diagonalisation]
Si :
\[
A=PDP^{-1},
\]
avec :
\[
D=\diag(\lambda_1,\dots,\lambda_n),
\]
alors la \(i\)-ième colonne de \(P\) est un vecteur propre de \(A\) associé à \(\lambda_i\).
\end{methodbox}

%=========================================================
\section{Critères fondamentaux de diagonalisation}

\subsection{Critère par somme directe des sous-espaces propres}

\begin{theorembox}[Critère géométrique]
Un endomorphisme \(u\in\Lcal(E)\) est diagonalisable si et seulement si :
\[
E=\bigoplus_{\lambda\in\Sp(u)}E_\lambda(u).
\]
\end{theorembox}

\begin{proof}
Supposons \(u\) diagonalisable.  
Il existe une base de \(E\) formée de vecteurs propres.  
En regroupant les vecteurs de cette base selon leur valeur propre, on obtient une base de chaque sous-espace propre.  
La réunion de ces bases forme une base de \(E\).  
Donc :
\[
E=\bigoplus_{\lambda\in\Sp(u)}E_\lambda(u).
\]

Réciproquement, supposons :
\[
E=\bigoplus_{\lambda\in\Sp(u)}E_\lambda(u).
\]
Choisissons une base de chaque \(E_\lambda(u)\).  
La réunion de ces bases est une base de \(E\), et tous ses vecteurs sont propres.  
Donc \(u\) est diagonalisable.
\end{proof}

\begin{corollaire}[Critère dimensionnel]
Un endomorphisme \(u\) est diagonalisable si et seulement si :
\[
\sum_{\lambda\in\Sp(u)}\dim(E_\lambda(u))=\dim(E).
\]
\end{corollaire}

\begin{proof}
Les sous-espaces propres associés à des valeurs propres distinctes sont en somme directe.  
Donc la dimension de leur somme est la somme de leurs dimensions.  
Cette somme vaut \(E\) si et seulement si sa dimension est \(\dim(E)\).
\end{proof}

\begin{retenirbox}
Diagonaliser revient à vérifier qu'il y a assez de vecteurs propres pour former une base.
\end{retenirbox}

\subsection{Critère des valeurs propres distinctes}

\begin{theorembox}
Soit \(E\) de dimension \(n\).  
Si \(u\) possède \(n\) valeurs propres distinctes dans \(\K\), alors \(u\) est diagonalisable.
\end{theorembox}

\begin{proof}
Choisissons un vecteur propre non nul associé à chacune des \(n\) valeurs propres.  
Des vecteurs propres associés à des valeurs propres deux à deux distinctes sont linéairement indépendants.  
On obtient donc une famille libre de \(n\) vecteurs dans un espace de dimension \(n\).  
C'est une base de \(E\).  
Cette base est formée de vecteurs propres, donc \(u\) est diagonalisable.
\end{proof}

\begin{erreurbox}
Avoir moins de \(n\) valeurs propres distinctes ne signifie pas que l'endomorphisme n'est pas diagonalisable.  
Exemple :
\[
A=
\begin{pmatrix}
2&0&0\\
0&3&0\\
0&0&3
\end{pmatrix}
\]
est diagonale, donc diagonalisable, bien qu'elle n'ait que deux valeurs propres distinctes.
\end{erreurbox}

\subsection{Critère par multiplicités}

\begin{definitionbox}[Multiplicités]
Soit \(\lambda\in\Sp(u)\).  
La multiplicité algébrique de \(\lambda\), notée \(m_\lambda\), est sa multiplicité comme racine de \(\chi_u\).  
La multiplicité géométrique de \(\lambda\) est :
\[
\dim(E_\lambda(u)).
\]
\end{definitionbox}

\begin{theorembox}[Inégalité fondamentale]
Pour toute valeur propre \(\lambda\),
\[
1\leq \dim(E_\lambda(u))\leq m_\lambda.
\]
\end{theorembox}

\begin{proof}
Comme \(\lambda\) est valeur propre :
\[
E_\lambda(u)\neq\{0\}.
\]
Donc :
\[
\dim(E_\lambda(u))\geq1.
\]

Posons :
\[
r=\dim(E_\lambda(u)).
\]
Choisissons une base :
\[
(e_1,\dots,e_r)
\]
de \(E_\lambda(u)\), complétée en une base de \(E\).  
Dans cette base, la matrice de \(u\) est de la forme :
\[
\begin{pmatrix}
\lambda I_r & *\\
0 & B
\end{pmatrix}.
\]
Donc :
\[
\chi_u(X)=(X-\lambda)^r\chi_B(X).
\]
Ainsi \((X-\lambda)^r\) divise \(\chi_u\), ce qui donne :
\[
r\leq m_\lambda.
\]
\end{proof}

\begin{theorembox}[Critère complet de diagonalisation]
Soit \(E\) de dimension \(n\).  
L'endomorphisme \(u\in\Lcal(E)\) est diagonalisable sur \(\K\) si et seulement si :
\begin{enumerate}
    \item \(\chi_u\) est scindé sur \(\K\) ;
    \item pour toute valeur propre \(\lambda\),
    \[
    \dim(E_\lambda(u))=m_\lambda.
    \]
\end{enumerate}
\end{theorembox}

\begin{proof}
Supposons \(u\) diagonalisable.  
Dans une base de vecteurs propres, sa matrice est diagonale :
\[
D=\diag(\lambda_1,\dots,\lambda_n).
\]
Donc :
\[
\chi_u(X)=\prod_{i=1}^n(X-\lambda_i),
\]
ce qui montre que \(\chi_u\) est scindé.  
Pour une valeur propre \(\lambda\), le nombre d'occurrences de \(\lambda\) sur la diagonale est à la fois sa multiplicité algébrique et la dimension de son sous-espace propre.  
Donc :
\[
\dim(E_\lambda)=m_\lambda.
\]

Réciproquement, supposons les deux conditions vérifiées.  
Comme \(\chi_u\) est scindé de degré \(n\),
\[
\sum_{\lambda\in\Sp(u)}m_\lambda=n.
\]
Par hypothèse :
\[
\sum_{\lambda\in\Sp(u)}\dim(E_\lambda)
=
\sum_{\lambda\in\Sp(u)}m_\lambda
=
n.
\]
Par le critère dimensionnel, \(u\) est diagonalisable.
\end{proof}

\begin{erreurbox}
Un polynôme caractéristique scindé ne suffit pas.  
Il faut encore vérifier :
\[
\dim(E_\lambda)=m_\lambda
\]
pour chaque valeur propre.
\end{erreurbox}

%=========================================================
\section{Polynômes annulateurs et diagonalisation}

\subsection{Polynôme annulateur}

\begin{definitionbox}[Polynôme annulateur]
Soit \(P\in\K[X]\).  
On dit que \(P\) est un polynôme annulateur de \(u\) si :
\[
P(u)=0.
\]
\end{definitionbox}

\begin{theorembox}
Si \(P(u)=0\), alors :
\[
\Sp(u)\subset\{\lambda\in\K:\ P(\lambda)=0\}.
\]
\end{theorembox}

\begin{proof}
Soit \(\lambda\in\Sp(u)\), et soit \(x\neq0\) tel que :
\[
u(x)=\lambda x.
\]
Alors, pour tout \(k\in\N\),
\[
u^k(x)=\lambda^kx.
\]
Si :
\[
P(X)=a_0+a_1X+\cdots+a_mX^m,
\]
alors :
\[
P(u)(x)=P(\lambda)x.
\]
Comme \(P(u)=0\), on a :
\[
P(\lambda)x=0.
\]
Puisque \(x\neq0\), il vient :
\[
P(\lambda)=0.
\]
\end{proof}

\subsection{Critère par polynôme annulateur}

\begin{theorembox}[Critère par polynôme annulateur scindé simple]
Si \(u\) admet un polynôme annulateur scindé à racines simples, alors \(u\) est diagonalisable.
\end{theorembox}

\begin{proof}
Supposons :
\[
P(X)=\prod_{i=1}^r(X-\lambda_i),
\]
avec les \(\lambda_i\) deux à deux distincts, et :
\[
P(u)=0.
\]
Les polynômes \(X-\lambda_i\) sont deux à deux premiers entre eux.  
Par le lemme de décomposition des noyaux :
\[
\Ker(P(u))
=
\bigoplus_{i=1}^r \Ker(u-\lambda_i\Id_E).
\]
Or :
\[
P(u)=0,
\]
donc :
\[
\Ker(P(u))=E.
\]
Ainsi :
\[
E=
\bigoplus_{i=1}^r E_{\lambda_i}(u).
\]
Donc \(u\) est diagonalisable.
\end{proof}

\begin{retenirbox}
Un polynôme annulateur scindé à racines simples suffit pour diagonaliser.
\[
P(u)=0,\quad P=\prod_i(X-\lambda_i),\quad \lambda_i\text{ distincts}
\Longrightarrow
u\text{ diagonalisable}.
\]
\end{retenirbox}

\begin{erreurbox}
Un polynôme annulateur scindé mais avec racines multiples ne garantit pas la diagonalisation.  
Exemple :
\[
A=
\begin{pmatrix}
1&1\\
0&1
\end{pmatrix}
\]
est annulée par \((X-1)^2\), mais n'est pas diagonalisable.
\end{erreurbox}

\subsection{Cas classiques}

\begin{theorembox}[Projecteurs]
Si \(p^2=p\), alors \(p\) est diagonalisable.
\end{theorembox}

\begin{proof}
On a :
\[
p^2-p=0.
\]
Donc \(p\) est annulé par :
\[
P(X)=X^2-X=X(X-1).
\]
Ce polynôme est scindé à racines simples.  
Donc \(p\) est diagonalisable.
\end{proof}

\begin{theorembox}[Involutions]
Si \(s^2=\Id_E\) et si \(\operatorname{car}(\K)\neq2\), alors \(s\) est diagonalisable.
\end{theorembox}

\begin{proof}
On a :
\[
s^2-\Id_E=0.
\]
Donc \(s\) est annulé par :
\[
P(X)=X^2-1=(X-1)(X+1).
\]
Si \(\operatorname{car}(\K)\neq2\), les racines \(1\) et \(-1\) sont distinctes.  
Donc \(P\) est scindé à racines simples, et \(s\) est diagonalisable.
\end{proof}

\begin{theorembox}[Nilpotents]
Si \(u^m=0\) pour un certain \(m\geq1\), alors la seule valeur propre possible de \(u\) est \(0\).  
Si \(u\) est diagonalisable, alors \(u=0\).
\end{theorembox}

\begin{proof}
Le polynôme :
\[
P(X)=X^m
\]
annule \(u\).  
Donc toute valeur propre de \(u\) est racine de \(P\), donc vaut \(0\).

Si \(u\) est diagonalisable, sa matrice dans une base de vecteurs propres est diagonale, avec uniquement des zéros sur la diagonale.  
Cette matrice est donc la matrice nulle.  
Ainsi :
\[
u=0.
\]
\end{proof}

%=========================================================
\section{Diagonalisation effective d'une matrice}

\begin{methodbox}[Méthode complète]
Pour diagonaliser une matrice \(A\in\M_n(\K)\) :
\begin{enumerate}
    \item Calculer :
    \[
    \chi_A(X)=\det(XI_n-A).
    \]
    \item Factoriser \(\chi_A\) sur le corps \(\K\).
    \item Déterminer les valeurs propres.
    \item Pour chaque valeur propre \(\lambda\), résoudre :
    \[
    (A-\lambda I_n)X=0.
    \]
    \item Comparer :
    \[
    \dim(E_\lambda)
    \quad\text{et}\quad
    m_\lambda.
    \]
    \item Si la somme des dimensions des sous-espaces propres vaut \(n\), choisir une base de chaque sous-espace propre.
    \item Former \(P\) avec ces vecteurs propres en colonnes.
    \item Former \(D\) en plaçant sur la diagonale les valeurs propres dans le même ordre que les colonnes de \(P\).
    \item Conclure :
    \[
    A=PDP^{-1}.
    \]
\end{enumerate}
\end{methodbox}

\begin{erreurbox}
L'ordre des colonnes de \(P\) doit correspondre à l'ordre des valeurs propres sur la diagonale de \(D\).  
Si on échange deux colonnes dans \(P\), il faut échanger les valeurs propres correspondantes dans \(D\).
\end{erreurbox}

\subsection{Exemple détaillé : deux valeurs propres distinctes}

Soit :
\[
A=
\begin{pmatrix}
1&1\\
0&2
\end{pmatrix}.
\]

On calcule :
\[
\chi_A(X)=
\det
\begin{pmatrix}
X-1&-1\\
0&X-2
\end{pmatrix}
=
(X-1)(X-2).
\]
Les valeurs propres sont :
\[
1
\quad\text{et}\quad
2.
\]
Elles sont distinctes, donc \(A\) est diagonalisable.

Pour \(\lambda=1\) :
\[
A-I_2=
\begin{pmatrix}
0&1\\
0&1
\end{pmatrix}.
\]
Le système donne :
\[
y=0.
\]
Donc :
\[
E_1=\Vect
\begin{pmatrix}
1\\0
\end{pmatrix}.
\]

Pour \(\lambda=2\) :
\[
A-2I_2=
\begin{pmatrix}
-1&1\\
0&0
\end{pmatrix}.
\]
Le système donne :
\[
y=x.
\]
Donc :
\[
E_2=\Vect
\begin{pmatrix}
1\\1
\end{pmatrix}.
\]

On prend :
\[
P=
\begin{pmatrix}
1&1\\
0&1
\end{pmatrix},
\qquad
D=
\begin{pmatrix}
1&0\\
0&2
\end{pmatrix}.
\]
Alors :
\[
A=PDP^{-1}.
\]

\subsection{Exemple détaillé : non diagonalisable}

Soit :
\[
A=
\begin{pmatrix}
1&1\\
0&1
\end{pmatrix}.
\]
On a :
\[
\chi_A(X)=(X-1)^2.
\]
L'unique valeur propre est \(1\), de multiplicité algébrique \(2\).

Calculons :
\[
A-I_2=
\begin{pmatrix}
0&1\\
0&0
\end{pmatrix}.
\]
Le système donne :
\[
y=0.
\]
Donc :
\[
E_1=
\Vect
\begin{pmatrix}
1\\0
\end{pmatrix}.
\]
Ainsi :
\[
\dim(E_1)=1<2.
\]
La matrice n'est pas diagonalisable.

\begin{retenirbox}
La matrice
\[
\begin{pmatrix}
1&1\\
0&1
\end{pmatrix}
\]
est l'exemple fondamental d'une matrice non diagonalisable mais déjà triangulaire.
\end{retenirbox}

\subsection{Exemple détaillé : multiplicité algébrique supérieure à \(1\) mais diagonalisable}

Soit :
\[
A=
\begin{pmatrix}
2&0&0\\
0&3&0\\
0&0&3
\end{pmatrix}.
\]
On a :
\[
\chi_A(X)=(X-2)(X-3)^2.
\]
Pour \(\lambda=2\),
\[
E_2=\Vect
\begin{pmatrix}
1\\0\\0
\end{pmatrix}.
\]
Pour \(\lambda=3\),
\[
E_3=\Vect
\left(
\begin{pmatrix}
0\\1\\0
\end{pmatrix},
\begin{pmatrix}
0\\0\\1
\end{pmatrix}
\right).
\]
Donc :
\[
\dim(E_2)+\dim(E_3)=1+2=3.
\]
La matrice est diagonalisable.

%=========================================================
\section{Applications de la diagonalisation}

\subsection{Calcul de puissances}

\begin{theorembox}
Si :
\[
A=PDP^{-1},
\]
alors pour tout \(m\in\N\),
\[
A^m=PD^mP^{-1}.
\]
\end{theorembox}

\begin{proof}
Pour \(m=0\), on a :
\[
A^0=I_n=P I_n P^{-1}=PD^0P^{-1}.
\]
Pour \(m\geq1\), on utilise l'associativité :
\[
A^m=(PDP^{-1})^m.
\]
Les facteurs \(P^{-1}P\) se simplifient successivement :
\[
A^m=PD^mP^{-1}.
\]
\end{proof}

\begin{theorembox}
Si :
\[
D=\diag(\lambda_1,\dots,\lambda_n),
\]
alors :
\[
D^m=\diag(\lambda_1^m,\dots,\lambda_n^m).
\]
\end{theorembox}

\begin{proof}
Le produit de matrices diagonales se calcule coefficient diagonal par coefficient diagonal.
\end{proof}

\begin{methodbox}[Calculer \(A^n\)]
Pour calculer \(A^n\) :
\begin{enumerate}
    \item diagonaliser \(A=PDP^{-1}\) ;
    \item calculer \(D^n\) ;
    \item revenir à :
    \[
    A^n=PD^nP^{-1}.
    \]
\end{enumerate}
\end{methodbox}

\subsection{Polynômes de matrices}

\begin{theorembox}
Si :
\[
A=PDP^{-1}
\]
et si \(Q\in\K[X]\), alors :
\[
Q(A)=PQ(D)P^{-1}.
\]
\end{theorembox}

\begin{proof}
Écrivons :
\[
Q(X)=a_0+a_1X+\cdots+a_mX^m.
\]
Alors :
\[
Q(A)=a_0I_n+a_1A+\cdots+a_mA^m.
\]
Comme :
\[
A^k=PD^kP^{-1},
\]
on obtient :
\[
Q(A)
=
P(a_0I_n+a_1D+\cdots+a_mD^m)P^{-1}
=
PQ(D)P^{-1}.
\]
\end{proof}

\begin{corollaire}
Si :
\[
D=\diag(\lambda_1,\dots,\lambda_n),
\]
alors :
\[
Q(D)=\diag(Q(\lambda_1),\dots,Q(\lambda_n)).
\]
\end{corollaire}

\subsection{Suites vectorielles}

Une suite définie par :
\[
X_{n+1}=AX_n
\]
vérifie :
\[
X_n=A^nX_0.
\]
Si \(A\) est diagonalisable, on obtient :
\[
X_n=PD^nP^{-1}X_0.
\]

\subsection{Exemple : suite de Fibonacci}

La suite de Fibonacci vérifie :
\[
F_{n+2}=F_{n+1}+F_n.
\]
Posons :
\[
X_n=
\begin{pmatrix}
F_{n+1}\\
F_n
\end{pmatrix}.
\]
Alors :
\[
X_{n+1}
=
\begin{pmatrix}
1&1\\
1&0
\end{pmatrix}
X_n.
\]
La matrice :
\[
A=
\begin{pmatrix}
1&1\\
1&0
\end{pmatrix}
\]
a pour polynôme caractéristique :
\[
\chi_A(X)=X^2-X-1.
\]
Ses valeurs propres sont :
\[
\varphi=\frac{1+\sqrt5}{2},
\qquad
\psi=\frac{1-\sqrt5}{2}.
\]
Elles sont distinctes, donc \(A\) est diagonalisable.  
On peut alors exprimer \(A^n\), puis obtenir la formule de Binet :
\[
F_n=\frac{\varphi^n-\psi^n}{\sqrt5}.
\]

%=========================================================
\section{Trigonalisation : définition et motivation}

\subsection{Pourquoi trigonaliser ?}

La diagonalisation est idéale, mais elle n'est pas toujours possible.  
La matrice :
\[
A=
\begin{pmatrix}
1&1\\
0&1
\end{pmatrix}
\]
n'est pas diagonalisable, car son unique sous-espace propre est de dimension \(1\).  
Pourtant, elle est déjà très simple : elle est triangulaire supérieure.

La trigonalisation est donc une réduction moins forte que la diagonalisation, mais plus générale.

\begin{retenirbox}
La diagonalisation cherche une base de vecteurs propres.  
La trigonalisation cherche une base dans laquelle chaque vecteur est envoyé dans l'espace engendré par lui-même et les vecteurs précédents.
\end{retenirbox}

\subsection{Matrices triangulaires}

\begin{definitionbox}[Matrice triangulaire supérieure]
Une matrice \(T=(t_{ij})\in\M_n(\K)\) est triangulaire supérieure si :
\[
i>j\Longrightarrow t_{ij}=0.
\]
Autrement dit, tous les coefficients strictement sous la diagonale sont nuls.
\end{definitionbox}

\begin{theorembox}
Si \(T\) est triangulaire supérieure, alors :
\[
\chi_T(X)=\prod_{i=1}^n(X-t_{ii}).
\]
En particulier, les valeurs propres de \(T\) sont ses coefficients diagonaux.
\end{theorembox}

\begin{proof}
La matrice \(XI_n-T\) est triangulaire supérieure de coefficients diagonaux :
\[
X-t_{11},\dots,X-t_{nn}.
\]
Son déterminant est donc :
\[
\det(XI_n-T)=\prod_{i=1}^n(X-t_{ii}).
\]
\end{proof}

\begin{theorembox}
Si \(T\) est triangulaire supérieure, alors :
\[
\det(T)=t_{11}\cdots t_{nn},
\qquad
\tr(T)=t_{11}+\cdots+t_{nn}.
\]
\end{theorembox}

\begin{proof}
Le déterminant d'une matrice triangulaire est le produit de ses coefficients diagonaux.  
La trace est, par définition, la somme des coefficients diagonaux.
\end{proof}

\subsection{Endomorphisme trigonalisable}

\begin{definitionbox}[Endomorphisme trigonalisable]
Un endomorphisme \(u\in\Lcal(E)\) est dit trigonalisable s'il existe une base de \(E\) dans laquelle la matrice de \(u\) est triangulaire supérieure.
\end{definitionbox}

\begin{definitionbox}[Matrice trigonalisable]
Une matrice \(A\in\M_n(\K)\) est trigonalisable s'il existe \(P\in GL_n(\K)\) et une matrice triangulaire supérieure \(T\) telles que :
\[
A=PTP^{-1}.
\]
\end{definitionbox}

\begin{retenirbox}
Toute matrice diagonalisable est trigonalisable, car toute matrice diagonale est triangulaire.  
La réciproque est fausse.
\end{retenirbox}

%=========================================================
\section{Sous-espaces stables et drapeaux}

\subsection{Sous-espace stable}

\begin{definitionbox}[Sous-espace stable]
Soit \(u\in\Lcal(E)\).  
Un sous-espace vectoriel \(F\subset E\) est stable par \(u\) si :
\[
u(F)\subset F.
\]
\end{definitionbox}

\begin{exercicebox}[Sous-espace propre stable]
Montrer que tout sous-espace propre \(E_\lambda(u)\) est stable par \(u\).
\end{exercicebox}

\begin{correctionbox}
Soit \(x\in E_\lambda(u)\).  
Alors :
\[
u(x)=\lambda x.
\]
Comme \(E_\lambda(u)\) est un sous-espace vectoriel et que \(x\in E_\lambda(u)\), on a :
\[
\lambda x\in E_\lambda(u).
\]
Donc :
\[
u(x)\in E_\lambda(u).
\]
Ainsi :
\[
u(E_\lambda(u))\subset E_\lambda(u).
\]
\end{correctionbox}

\subsection{Matrice par blocs associée à un sous-espace stable}

\begin{theorembox}
Soit \(F\) un sous-espace stable par \(u\).  
Si l'on choisit une base de \(F\), complétée en une base de \(E\), alors la matrice de \(u\) dans cette base est de la forme :
\[
\begin{pmatrix}
A&B\\
0&C
\end{pmatrix}.
\]
\end{theorembox}

\begin{proof}
Soit :
\[
(e_1,\dots,e_p)
\]
une base de \(F\), complétée en :
\[
(e_1,\dots,e_p,e_{p+1},\dots,e_n)
\]
base de \(E\).

Comme \(F\) est stable, pour tout \(j\leq p\),
\[
u(e_j)\in F.
\]
Donc les coordonnées de \(u(e_j)\) selon les vecteurs \(e_{p+1},\dots,e_n\) sont nulles.  
Cela impose que le bloc inférieur gauche de la matrice soit nul.  
La matrice a donc la forme :
\[
\begin{pmatrix}
A&B\\
0&C
\end{pmatrix}.
\]
\end{proof}

\begin{theorembox}[Polynôme caractéristique d'une matrice triangulaire par blocs]
Si :
\[
M=
\begin{pmatrix}
A&B\\
0&C
\end{pmatrix},
\]
alors :
\[
\chi_M(X)=\chi_A(X)\chi_C(X).
\]
\end{theorembox}

\begin{proof}
On calcule :
\[
XI-M=
\begin{pmatrix}
XI-A&-B\\
0&XI-C
\end{pmatrix}.
\]
Cette matrice est triangulaire par blocs.  
Son déterminant est donc :
\[
\det(XI-M)=\det(XI-A)\det(XI-C).
\]
Ainsi :
\[
\chi_M=\chi_A\chi_C.
\]
\end{proof}

\subsection{Drapeaux stables}

\begin{definitionbox}[Drapeau]
Un drapeau de \(E\) est une suite de sous-espaces :
\[
\{0\}=E_0\subset E_1\subset E_2\subset\cdots\subset E_n=E
\]
telle que :
\[
\dim(E_k)=k.
\]
\end{definitionbox}

\begin{definitionbox}[Drapeau stable]
Un drapeau est stable par \(u\) si chaque \(E_k\) est stable par \(u\).
\end{definitionbox}

\begin{theorembox}[Caractérisation par drapeau stable]
Un endomorphisme \(u\) est trigonalisable si et seulement s'il existe un drapeau stable par \(u\).
\end{theorembox}

\begin{proof}
Supposons \(u\) trigonalisable.  
Il existe une base :
\[
(e_1,\dots,e_n)
\]
dans laquelle la matrice de \(u\) est triangulaire supérieure.  
Posons :
\[
E_k=\Vect(e_1,\dots,e_k).
\]
Comme la matrice est triangulaire supérieure, pour tout \(j\leq k\),
\[
u(e_j)\in\Vect(e_1,\dots,e_j)\subset E_k.
\]
Donc chaque \(E_k\) est stable.  
On obtient un drapeau stable.

Réciproquement, supposons qu'il existe un drapeau stable :
\[
\{0\}=E_0\subset E_1\subset\cdots\subset E_n=E.
\]
Pour chaque \(k\), choisissons :
\[
e_k\in E_k\setminus E_{k-1}.
\]
Alors \((e_1,\dots,e_n)\) est une base de \(E\).  
Comme \(E_k\) est stable, on a :
\[
u(e_k)\in E_k=\Vect(e_1,\dots,e_k).
\]
Cela signifie que la matrice de \(u\) dans cette base est triangulaire supérieure.
\end{proof}

%=========================================================
\section{Critère fondamental de trigonalisation}

\begin{theorembox}[Critère de trigonalisation]
Un endomorphisme \(u\in\Lcal(E)\), avec \(E\) de dimension finie, est trigonalisable sur \(\K\) si et seulement si son polynôme caractéristique \(\chi_u\) est scindé sur \(\K\).
\end{theorembox}

\subsection*{Sens direct}

\begin{proof}
Supposons \(u\) trigonalisable.  
Dans une base adaptée, sa matrice est triangulaire supérieure :
\[
T=
\begin{pmatrix}
\lambda_1&*&\cdots&*\\
0&\lambda_2&\cdots&*\\
\vdots&\vdots&\ddots&\vdots\\
0&0&\cdots&\lambda_n
\end{pmatrix}.
\]
Alors :
\[
\chi_u(X)=\chi_T(X)=\prod_{i=1}^n(X-\lambda_i).
\]
Donc \(\chi_u\) est scindé sur \(\K\).
\end{proof}

\subsection*{Sens réciproque}

\begin{proof}
On raisonne par récurrence sur :
\[
n=\dim(E).
\]

Si \(n=1\), tout endomorphisme est représenté par une matrice \(1\times1\), donc est trigonalisable.

Supposons le résultat vrai en dimension \(n-1\).  
Soit \(E\) de dimension \(n\), et soit \(u\in\Lcal(E)\) tel que \(\chi_u\) soit scindé.

Comme \(\chi_u\) est scindé et de degré \(n\), il admet une racine :
\[
\lambda\in\K.
\]
Donc \(\lambda\) est une valeur propre de \(u\).  
Il existe alors \(e_1\neq0\) tel que :
\[
u(e_1)=\lambda e_1.
\]
La droite :
\[
E_1=\Vect(e_1)
\]
est stable par \(u\).

Complétons \(e_1\) en une base de \(E\).  
Dans cette base, la matrice de \(u\) est de la forme :
\[
\begin{pmatrix}
\lambda&*\\
0&B
\end{pmatrix}.
\]
Le polynôme caractéristique se factorise alors :
\[
\chi_u(X)=(X-\lambda)\chi_B(X).
\]
Comme \(\chi_u\) est scindé, \(\chi_B\) est scindé.

Par hypothèse de récurrence, la matrice \(B\) est trigonalisable.  
Il existe donc une base du complément quotient dans laquelle \(B\) devient triangulaire supérieure.  
En relevant cette base dans \(E\) et en la précédant par \(e_1\), on obtient une base de \(E\) dans laquelle la matrice de \(u\) est triangulaire supérieure.

Donc \(u\) est trigonalisable.
\end{proof}

\begin{retenirbox}
Sur \(\C\), tout polynôme est scindé.  
Donc toute matrice complexe est trigonalisable.
\end{retenirbox}

\begin{corollaire}
Toute matrice de \(\M_n(\C)\) est trigonalisable sur \(\C\).
\end{corollaire}

\begin{proof}
Le polynôme caractéristique d'une matrice complexe est un polynôme de \(\C[X]\).  
Par le théorème fondamental de l'algèbre, il est scindé sur \(\C\).  
La matrice est donc trigonalisable.
\end{proof}

\begin{erreurbox}
Sur \(\R\), toute matrice n'est pas trigonalisable.  
Exemple :
\[
A=
\begin{pmatrix}
0&-1\\
1&0
\end{pmatrix}
\]
a pour polynôme caractéristique :
\[
X^2+1,
\]
qui n'est pas scindé sur \(\R\).
\end{erreurbox}

%=========================================================
\section{Diagonalisation versus trigonalisation}

\begin{theorembox}
Toute matrice diagonalisable est trigonalisable.
\end{theorembox}

\begin{proof}
Si \(A\) est diagonalisable, il existe :
\[
A=PDP^{-1}
\]
avec \(D\) diagonale.  
Or toute matrice diagonale est triangulaire supérieure.  
Donc \(A\) est trigonalisable.
\end{proof}

\begin{theorembox}
La réciproque est fausse.
\end{theorembox}

\begin{proof}
La matrice :
\[
A=
\begin{pmatrix}
1&1\\
0&1
\end{pmatrix}
\]
est triangulaire supérieure, donc trigonalisable.  
Mais :
\[
\chi_A(X)=(X-1)^2,
\]
et :
\[
E_1(A)=\Ker(A-I_2)=\Vect
\begin{pmatrix}
1\\0
\end{pmatrix}.
\]
Donc :
\[
\dim(E_1)=1<2.
\]
Elle n'est pas diagonalisable.
\end{proof}

\begin{center}
\begin{tabularx}{0.95\linewidth}{|X|X|}
\hline
\textbf{Diagonalisation} & \textbf{Trigonalisation}\\
\hline
Base de vecteurs propres & Base adaptée à un drapeau stable\\
\hline
Matrice diagonale & Matrice triangulaire supérieure\\
\hline
Très forte réduction & Réduction plus faible\\
\hline
Pas toujours possible, même si \(\chi\) est scindé & Possible dès que \(\chi\) est scindé\\
\hline
Permet de calculer très simplement \(A^n\) & Permet de lire trace, déterminant, valeurs propres\\
\hline
\end{tabularx}
\end{center}

%=========================================================
\section{Ce que donne une matrice triangulaire}

\subsection{Lecture des valeurs propres}

Si :
\[
A=PTP^{-1}
\]
avec \(T\) triangulaire supérieure :
\[
T=
\begin{pmatrix}
\lambda_1&*&\cdots&*\\
0&\lambda_2&\cdots&*\\
\vdots&\vdots&\ddots&\vdots\\
0&0&\cdots&\lambda_n
\end{pmatrix},
\]
alors :
\[
\chi_A(X)=\prod_{i=1}^n(X-\lambda_i).
\]
Les valeurs propres de \(A\), comptées avec multiplicité, sont donc :
\[
\lambda_1,\dots,\lambda_n.
\]

\subsection{Trace et déterminant}

\begin{theorembox}
Si \(A\) est trigonalisable et si ses valeurs propres comptées avec multiplicité sont :
\[
\lambda_1,\dots,\lambda_n,
\]
alors :
\[
\tr(A)=\lambda_1+\cdots+\lambda_n,
\]
et :
\[
\det(A)=\lambda_1\cdots\lambda_n.
\]
\end{theorembox}

\begin{proof}
Soit \(T\) triangulaire supérieure semblable à \(A\), avec diagonale :
\[
\lambda_1,\dots,\lambda_n.
\]
La trace et le déterminant sont invariants par similitude.  
Donc :
\[
\tr(A)=\tr(T)=\lambda_1+\cdots+\lambda_n,
\]
et :
\[
\det(A)=\det(T)=\lambda_1\cdots\lambda_n.
\]
\end{proof}

\subsection{Polynômes d'une matrice triangulaire}

\begin{theorembox}
Si \(T\) est triangulaire supérieure et si \(P\in\K[X]\), alors \(P(T)\) est triangulaire supérieure.
\end{theorembox}

\begin{proof}
Le produit de deux matrices triangulaires supérieures est triangulaire supérieur.  
Donc toutes les puissances :
\[
T^k
\]
sont triangulaires supérieures.  
Une combinaison linéaire de matrices triangulaires supérieures est encore triangulaire supérieure.  
Ainsi :
\[
P(T)
\]
est triangulaire supérieure.
\end{proof}

\begin{theorembox}
Si \(T\) est triangulaire supérieure de diagonale :
\[
\lambda_1,\dots,\lambda_n,
\]
alors \(P(T)\) est triangulaire supérieure de diagonale :
\[
P(\lambda_1),\dots,P(\lambda_n).
\]
\end{theorembox}

\begin{proof}
La diagonale de \(T^k\) est :
\[
\lambda_1^k,\dots,\lambda_n^k.
\]
Si :
\[
P(X)=a_0+a_1X+\cdots+a_mX^m,
\]
alors la diagonale de \(P(T)\) est :
\[
a_0+a_1\lambda_i+\cdots+a_m\lambda_i^m=P(\lambda_i).
\]
\end{proof}

%=========================================================
\section{Sous-espaces caractéristiques et idée de Jordan}

\subsection{Sous-espaces caractéristiques}

\begin{definitionbox}[Sous-espace caractéristique]
Soit \(\lambda\in\Sp(u)\).  
Pour \(m\geq1\), on considère :
\[
\Ker((u-\lambda\Id_E)^m).
\]
Lorsque \(m\) est assez grand, ces noyaux décrivent les vecteurs qui deviennent propres après application répétée de \(u-\lambda\Id_E\).
\end{definitionbox}

\begin{retenirbox}
Quand \(u\) n'est pas diagonalisable, les sous-espaces propres peuvent être trop petits.  
Les noyaux :
\[
\Ker((u-\lambda\Id_E)^m)
\]
permettent de comprendre les directions propres généralisées.
\end{retenirbox}

\begin{erreurbox}
Dans un premier cours de prépa, il ne faut pas confondre :
\[
\Ker(u-\lambda\Id_E)
\]
et :
\[
\Ker((u-\lambda\Id_E)^m).
\]
Le premier est le sous-espace propre.  
Le second contient des vecteurs propres généralisés.
\end{erreurbox}

\subsection{Bloc de Jordan élémentaire}

La matrice :
\[
J_\lambda=
\begin{pmatrix}
\lambda&1\\
0&\lambda
\end{pmatrix}
\]
est appelée un bloc de Jordan de taille \(2\).  
Elle n'est pas diagonalisable, mais elle est triangulaire supérieure.

On peut écrire :
\[
J_\lambda=\lambda I_2+N,
\]
où :
\[
N=
\begin{pmatrix}
0&1\\
0&0
\end{pmatrix},
\qquad
N^2=0.
\]

\begin{retenirbox}
La partie diagonale donne la valeur propre.  
La partie nilpotente mesure l'obstruction à la diagonalisation.
\end{retenirbox}

%=========================================================
\section{Méthodes pratiques de trigonalisation}

\begin{methodbox}[Vérifier la trigonalisabilité]
Pour savoir si \(A\in\M_n(\K)\) est trigonalisable :
\begin{enumerate}
    \item calculer \(\chi_A\) ;
    \item vérifier si \(\chi_A\) est scindé sur \(\K\) ;
    \item si oui, \(A\) est trigonalisable ;
    \item si non, \(A\) ne l'est pas sur \(\K\).
\end{enumerate}
\end{methodbox}

\begin{methodbox}[Construire une trigonalisation en petite dimension]
Pour construire une base triangulante :
\begin{enumerate}
    \item choisir un vecteur propre \(e_1\) ;
    \item compléter en une base ;
    \item obtenir une matrice triangulaire par blocs ;
    \item recommencer sur le bloc quotient ;
    \item construire progressivement une base :
    \[
    (e_1,\dots,e_n)
    \]
    telle que :
    \[
    u(e_k)\in\Vect(e_1,\dots,e_k).
    \]
\end{enumerate}
\end{methodbox}

\begin{erreurbox}
Le programme demande surtout de savoir caractériser la trigonalisabilité et exploiter une forme triangulaire.  
La construction effective d'une base trigonalante peut être délicate et n'est pas toujours l'objectif prioritaire.
\end{erreurbox}

%=========================================================
\section{Exemples détaillés}

\subsection{Exemple 1 : trigonalisable mais non diagonalisable}

Soit :
\[
A=
\begin{pmatrix}
1&1\\
0&1
\end{pmatrix}.
\]
La matrice est déjà triangulaire supérieure.  
Donc elle est trigonalisable.

Son polynôme caractéristique est :
\[
\chi_A(X)=(X-1)^2.
\]
Elle est scindée sur tout corps contenant \(1\).

Son sous-espace propre associé à \(1\) est :
\[
E_1=\Vect
\begin{pmatrix}
1\\0
\end{pmatrix}.
\]
Donc :
\[
\dim(E_1)=1<2.
\]
Elle n'est pas diagonalisable.

\subsection{Exemple 2 : non trigonalisable sur \(\R\), trigonalisable sur \(\C\)}

Soit :
\[
A=
\begin{pmatrix}
0&-1\\
1&0
\end{pmatrix}.
\]
On calcule :
\[
\chi_A(X)=X^2+1.
\]
Sur \(\R\), ce polynôme n'est pas scindé.  
Donc \(A\) n'est pas trigonalisable sur \(\R\).

Sur \(\C\), on a :
\[
X^2+1=(X-i)(X+i).
\]
Donc \(A\) est trigonalisable sur \(\C\), et même diagonalisable car les valeurs propres sont distinctes.

\subsection{Exemple 3 : triangulaire \(3\times3\)}

Soit :
\[
A=
\begin{pmatrix}
2&1&0\\
0&2&1\\
0&0&3
\end{pmatrix}.
\]
La matrice est triangulaire supérieure, donc trigonalisable.  
Son polynôme caractéristique est :
\[
\chi_A(X)=(X-2)^2(X-3).
\]
Les valeurs propres sont :
\[
2
\quad\text{et}\quad
3.
\]

Pour \(\lambda=2\),
\[
A-2I_3=
\begin{pmatrix}
0&1&0\\
0&0&1\\
0&0&1
\end{pmatrix}.
\]
Le système donne :
\[
y=0,\qquad z=0.
\]
Donc :
\[
E_2=\Vect
\begin{pmatrix}
1\\0\\0
\end{pmatrix}.
\]
Ainsi :
\[
\dim(E_2)=1<2.
\]
La matrice n'est pas diagonalisable.

%=========================================================
\newpage
\section{Exercices progressifs}

\subsection*{Thème 1 -- Diagonalisation}

\begin{exercicebox}[1]
Diagonaliser, si possible :
\[
A=
\begin{pmatrix}
1&1\\
0&2
\end{pmatrix}.
\]
\end{exercicebox}

\begin{exercicebox}[2]
Étudier la diagonalisabilité de :
\[
A=
\begin{pmatrix}
1&1\\
0&1
\end{pmatrix}.
\]
\end{exercicebox}

\begin{exercicebox}[3]
Diagonaliser :
\[
A=
\begin{pmatrix}
2&0&0\\
0&3&0\\
0&0&3
\end{pmatrix}.
\]
\end{exercicebox}

\begin{exercicebox}[4]
Soit :
\[
A=
\begin{pmatrix}
3&1&0\\
0&3&0\\
0&0&5
\end{pmatrix}.
\]
Étudier la diagonalisabilité de \(A\).
\end{exercicebox}

\subsection*{Thème 2 -- Polynômes annulateurs}

\begin{exercicebox}[5]
Soit \(p\in\Lcal(E)\) tel que :
\[
p^2=p.
\]
Montrer que \(p\) est diagonalisable et donner ses valeurs propres possibles.
\end{exercicebox}

\begin{exercicebox}[6]
Soit \(s\in\Lcal(E)\) tel que :
\[
s^2=\Id_E.
\]
Montrer que \(s\) est diagonalisable si \(\operatorname{car}(\K)\neq2\).
\end{exercicebox}

\begin{exercicebox}[7]
Soit \(u\in\Lcal(E)\) tel que :
\[
u^3-u=0.
\]
Montrer que \(u\) est diagonalisable si \(\operatorname{car}(\K)\neq2\).
\end{exercicebox}

\subsection*{Thème 3 -- Trigonalisation}

\begin{exercicebox}[8]
Montrer que toute matrice triangulaire supérieure est trigonalisable.
\end{exercicebox}

\begin{exercicebox}[9]
Étudier la trigonalisabilité sur \(\R\) et sur \(\C\) de :
\[
A=
\begin{pmatrix}
0&-1\\
1&0
\end{pmatrix}.
\]
\end{exercicebox}

\begin{exercicebox}[10]
Soit \(A\in\M_n(\C)\).  
Montrer que \(A\) est trigonalisable sur \(\C\).
\end{exercicebox}

\subsection*{Thème 4 -- Applications}

\begin{exercicebox}[11]
Soit :
\[
A=
\begin{pmatrix}
1&1\\
1&1
\end{pmatrix}.
\]
Diagonaliser \(A\), puis calculer \(A^n\) pour \(n\geq1\).
\end{exercicebox}

\begin{exercicebox}[12]
Soit \(A\) trigonalisable de valeurs propres comptées avec multiplicité :
\[
\lambda_1,\dots,\lambda_n.
\]
Montrer que :
\[
\tr(A)=\sum_{i=1}^n\lambda_i,
\qquad
\det(A)=\prod_{i=1}^n\lambda_i.
\]
\end{exercicebox}

%=========================================================
\newpage
\section{Corrigés détaillés des exercices}

\begin{correctionbox}[1]
On a :
\[
A=
\begin{pmatrix}
1&1\\
0&2
\end{pmatrix}.
\]
Le polynôme caractéristique est :
\[
\chi_A(X)=(X-1)(X-2).
\]
Les valeurs propres sont distinctes, donc \(A\) est diagonalisable.

Pour \(\lambda=1\),
\[
A-I=
\begin{pmatrix}
0&1\\
0&1
\end{pmatrix}.
\]
Donc :
\[
E_1=\Vect
\begin{pmatrix}
1\\0
\end{pmatrix}.
\]
Pour \(\lambda=2\),
\[
A-2I=
\begin{pmatrix}
-1&1\\
0&0
\end{pmatrix}.
\]
Donc :
\[
E_2=\Vect
\begin{pmatrix}
1\\1
\end{pmatrix}.
\]
On pose :
\[
P=
\begin{pmatrix}
1&1\\
0&1
\end{pmatrix},
\qquad
D=
\begin{pmatrix}
1&0\\
0&2
\end{pmatrix}.
\]
Alors :
\[
A=PDP^{-1}.
\]
\end{correctionbox}

\begin{correctionbox}[2]
On a :
\[
\chi_A(X)=(X-1)^2.
\]
L'unique valeur propre est \(1\).  
Or :
\[
A-I=
\begin{pmatrix}
0&1\\
0&0
\end{pmatrix}.
\]
Donc :
\[
E_1=\Vect
\begin{pmatrix}
1\\0
\end{pmatrix}.
\]
Ainsi :
\[
\dim(E_1)=1<2.
\]
La matrice n'est pas diagonalisable.
\end{correctionbox}

\begin{correctionbox}[3]
La matrice est déjà diagonale :
\[
A=\diag(2,3,3).
\]
Elle est donc diagonalisable.  
On peut prendre :
\[
P=I_3,
\qquad
D=A.
\]
Les sous-espaces propres sont :
\[
E_2=\Vect
\begin{pmatrix}
1\\0\\0
\end{pmatrix},
\]
et :
\[
E_3=\Vect
\left(
\begin{pmatrix}
0\\1\\0
\end{pmatrix},
\begin{pmatrix}
0\\0\\1
\end{pmatrix}
\right).
\]
\end{correctionbox}

\begin{correctionbox}[4]
La matrice est triangulaire supérieure.  
Donc :
\[
\chi_A(X)=(X-3)^2(X-5).
\]
Pour \(\lambda=3\),
\[
A-3I=
\begin{pmatrix}
0&1&0\\
0&0&0\\
0&0&2
\end{pmatrix}.
\]
Le système donne :
\[
y=0,\qquad z=0.
\]
Donc :
\[
E_3=\Vect
\begin{pmatrix}
1\\0\\0
\end{pmatrix}.
\]
Ainsi :
\[
\dim(E_3)=1<2.
\]
La matrice n'est pas diagonalisable.
\end{correctionbox}

\begin{correctionbox}[5]
Comme :
\[
p^2=p,
\]
on a :
\[
p^2-p=0.
\]
Donc \(p\) est annulé par :
\[
P(X)=X^2-X=X(X-1).
\]
Ce polynôme est scindé à racines simples.  
Donc \(p\) est diagonalisable.  
Ses valeurs propres possibles sont :
\[
0
\quad\text{et}\quad
1.
\]
\end{correctionbox}

\begin{correctionbox}[6]
Comme :
\[
s^2=\Id_E,
\]
on a :
\[
s^2-\Id_E=0.
\]
Donc \(s\) est annulé par :
\[
P(X)=X^2-1=(X-1)(X+1).
\]
Si \(\operatorname{car}(\K)\neq2\), les racines \(1\) et \(-1\) sont distinctes.  
Donc \(P\) est scindé à racines simples.  
Ainsi \(s\) est diagonalisable.
\end{correctionbox}

\begin{correctionbox}[7]
On a :
\[
u^3-u=0.
\]
Donc \(u\) est annulé par :
\[
P(X)=X^3-X=X(X-1)(X+1).
\]
Si \(\operatorname{car}(\K)\neq2\), les racines \(-1,0,1\) sont distinctes.  
Le polynôme est scindé à racines simples.  
Donc \(u\) est diagonalisable.
\end{correctionbox}

\begin{correctionbox}[8]
Si \(A\) est triangulaire supérieure, alors elle est semblable à une matrice triangulaire supérieure, elle-même, avec :
\[
P=I_n.
\]
Donc \(A\) est trigonalisable.
\end{correctionbox}

\begin{correctionbox}[9]
On calcule :
\[
\chi_A(X)=X^2+1.
\]
Sur \(\R\), ce polynôme n'est pas scindé.  
Donc \(A\) n'est pas trigonalisable sur \(\R\).

Sur \(\C\), on a :
\[
X^2+1=(X-i)(X+i).
\]
Le polynôme est scindé.  
Donc \(A\) est trigonalisable sur \(\C\).  
Comme les valeurs propres \(i\) et \(-i\) sont distinctes, elle est même diagonalisable sur \(\C\).
\end{correctionbox}

\begin{correctionbox}[10]
Le polynôme caractéristique de \(A\) est un polynôme de \(\C[X]\).  
Par le théorème fondamental de l'algèbre, il est scindé sur \(\C\).  
Par le critère fondamental de trigonalisation, \(A\) est trigonalisable sur \(\C\).
\end{correctionbox}

\begin{correctionbox}[11]
On calcule :
\[
\chi_A(X)=
\det
\begin{pmatrix}
X-1&-1\\
-1&X-1
\end{pmatrix}
=
(X-1)^2-1
=
X^2-2X
=
X(X-2).
\]
Les valeurs propres sont :
\[
0
\quad\text{et}\quad
2.
\]
Elles sont distinctes, donc \(A\) est diagonalisable.

Pour \(\lambda=0\),
\[
E_0=\Vect
\begin{pmatrix}
1\\-1
\end{pmatrix}.
\]
Pour \(\lambda=2\),
\[
E_2=\Vect
\begin{pmatrix}
1\\1
\end{pmatrix}.
\]
On prend :
\[
P=
\begin{pmatrix}
1&1\\
-1&1
\end{pmatrix},
\qquad
D=
\begin{pmatrix}
0&0\\
0&2
\end{pmatrix}.
\]
Alors :
\[
A=PDP^{-1}.
\]
Comme :
\[
P^{-1}=\frac12
\begin{pmatrix}
1&-1\\
1&1
\end{pmatrix},
\]
on obtient, pour \(n\geq1\),
\[
A^n=PD^nP^{-1}
=
2^{n-1}
\begin{pmatrix}
1&1\\
1&1
\end{pmatrix}.
\]
\end{correctionbox}

\begin{correctionbox}[12]
Comme \(A\) est trigonalisable, il existe \(P\in GL_n(\K)\) et \(T\) triangulaire supérieure tels que :
\[
A=PTP^{-1}.
\]
La diagonale de \(T\) est formée des valeurs propres de \(A\), comptées avec multiplicité :
\[
\lambda_1,\dots,\lambda_n.
\]
La trace et le déterminant sont invariants par similitude.  
Donc :
\[
\tr(A)=\tr(T)=\lambda_1+\cdots+\lambda_n,
\]
et :
\[
\det(A)=\det(T)=\lambda_1\cdots\lambda_n.
\]
\end{correctionbox}

%=========================================================
\newpage
\section{Grand devoir bilan final}

\begin{center}
\Large\textbf{Devoir bilan -- Diagonalisation et trigonalisation}
\end{center}

\subsection*{Exercice 1 -- Questions de cours}

Soit \(E\) un \(\K\)-espace vectoriel de dimension finie et \(u\in\Lcal(E)\).

\begin{enumerate}
    \item Définir un endomorphisme diagonalisable.
    \item Définir une matrice diagonalisable.
    \item Montrer que \(u\) est diagonalisable si et seulement si :
    \[
    E=\bigoplus_{\lambda\in\Sp(u)}E_\lambda(u).
    \]
    \item Énoncer et démontrer le critère des \(n\) valeurs propres distinctes.
    \item Définir un endomorphisme trigonalisable.
    \item Énoncer le critère fondamental de trigonalisation.
\end{enumerate}

\subsection*{Exercice 2 -- Diagonalisation complète}

Soit :
\[
A=
\begin{pmatrix}
2&1&0\\
0&2&0\\
0&0&3
\end{pmatrix}.
\]

\begin{enumerate}
    \item Calculer \(\chi_A\).
    \item Déterminer les valeurs propres.
    \item Calculer les sous-espaces propres.
    \item Donner les multiplicités algébriques et géométriques.
    \item La matrice est-elle diagonalisable ?
    \item Est-elle trigonalisable ?
\end{enumerate}

\subsection*{Exercice 3 -- Polynôme annulateur}

Soit \(u\in\Lcal(E)\) tel que :
\[
u^3-u=0.
\]

\begin{enumerate}
    \item Factoriser le polynôme \(X^3-X\).
    \item Déterminer les valeurs propres possibles de \(u\).
    \item Montrer que \(u\) est diagonalisable si \(\operatorname{car}(\K)\neq2\).
\end{enumerate}

\subsection*{Exercice 4 -- Trigonalisation sur \(\R\) et sur \(\C\)}

Soit :
\[
A=
\begin{pmatrix}
0&-1\\
1&0
\end{pmatrix}.
\]

\begin{enumerate}
    \item Calculer \(\chi_A\).
    \item Étudier la diagonalisation sur \(\R\).
    \item Étudier la trigonalisation sur \(\R\).
    \item Étudier la diagonalisation et la trigonalisation sur \(\C\).
\end{enumerate}

\subsection*{Exercice 5 -- Puissances}

Soit :
\[
A=
\begin{pmatrix}
1&1\\
1&1
\end{pmatrix}.
\]

\begin{enumerate}
    \item Diagonaliser \(A\).
    \item Calculer \(A^n\) pour \(n\geq1\).
    \item En déduire le comportement de \(A^nX_0\) selon \(X_0\in\R^2\).
\end{enumerate}

\newpage

%=========================================================
\section{Correction détaillée du devoir bilan}

\subsection*{Correction de l'exercice 1}

\textbf{1.}  
Un endomorphisme \(u\in\Lcal(E)\) est diagonalisable s'il existe une base de \(E\) formée de vecteurs propres de \(u\).

\textbf{2.}  
Une matrice \(A\in\M_n(\K)\) est diagonalisable s'il existe \(P\in GL_n(\K)\) et une matrice diagonale \(D\) telles que :
\[
A=PDP^{-1}.
\]

\textbf{3.}  
Si \(u\) est diagonalisable, une base de vecteurs propres se regroupe selon les valeurs propres.  
Cela donne :
\[
E=\bigoplus_{\lambda\in\Sp(u)}E_\lambda(u).
\]
Réciproquement, si cette somme directe vaut \(E\), on choisit une base de chaque sous-espace propre ; leur réunion donne une base de \(E\) formée de vecteurs propres.  
Donc \(u\) est diagonalisable.

\textbf{4.}  
Si \(\dim(E)=n\) et si \(u\) possède \(n\) valeurs propres distinctes, on choisit un vecteur propre non nul pour chacune.  
Ces \(n\) vecteurs sont libres car associés à des valeurs propres distinctes.  
Une famille libre de \(n\) vecteurs dans un espace de dimension \(n\) est une base.  
Donc \(u\) est diagonalisable.

\textbf{5.}  
Un endomorphisme \(u\) est trigonalisable s'il existe une base dans laquelle sa matrice est triangulaire supérieure.

\textbf{6.}  
Le critère fondamental est :
\[
u\text{ trigonalisable sur }\K
\Longleftrightarrow
\chi_u\text{ est scindé sur }\K.
\]

\subsection*{Correction de l'exercice 2}

La matrice :
\[
A=
\begin{pmatrix}
2&1&0\\
0&2&0\\
0&0&3
\end{pmatrix}
\]
est triangulaire supérieure.  
Donc :
\[
\chi_A(X)=(X-2)^2(X-3).
\]
Les valeurs propres sont :
\[
2
\quad\text{et}\quad
3.
\]

Pour \(\lambda=2\),
\[
A-2I_3=
\begin{pmatrix}
0&1&0\\
0&0&0\\
0&0&1
\end{pmatrix}.
\]
Le système donne :
\[
y=0,\qquad z=0.
\]
Donc :
\[
E_2=\Vect
\begin{pmatrix}
1\\0\\0
\end{pmatrix}.
\]
Ainsi :
\[
m_2=2,
\qquad
\dim(E_2)=1.
\]

Pour \(\lambda=3\),
\[
A-3I_3=
\begin{pmatrix}
-1&1&0\\
0&-1&0\\
0&0&0
\end{pmatrix}.
\]
Le système donne :
\[
x=0,\qquad y=0.
\]
Donc :
\[
E_3=\Vect
\begin{pmatrix}
0\\0\\1
\end{pmatrix}.
\]
Ainsi :
\[
m_3=1,
\qquad
\dim(E_3)=1.
\]

La somme des dimensions des sous-espaces propres vaut :
\[
1+1=2<3.
\]
La matrice n'est donc pas diagonalisable.

En revanche, elle est déjà triangulaire supérieure.  
Donc elle est trigonalisable.

\subsection*{Correction de l'exercice 3}

On factorise :
\[
X^3-X=X(X^2-1)=X(X-1)(X+1).
\]
Les valeurs propres possibles de \(u\) sont donc :
\[
-1,\quad0,\quad1.
\]
Si \(\operatorname{car}(\K)\neq2\), ces trois racines sont distinctes.  
Le polynôme annulateur :
\[
X(X-1)(X+1)
\]
est donc scindé à racines simples.  
Par le critère de diagonalisation par polynôme annulateur, \(u\) est diagonalisable.

\subsection*{Correction de l'exercice 4}

On calcule :
\[
\chi_A(X)=
\det
\begin{pmatrix}
X&1\\
-1&X
\end{pmatrix}
=
X^2+1.
\]

Sur \(\R\), ce polynôme n'a aucune racine.  
Donc \(A\) n'a aucune valeur propre réelle.  
Elle n'est donc pas diagonalisable sur \(\R\).  
Comme \(\chi_A\) n'est pas scindé sur \(\R\), elle n'est pas trigonalisable sur \(\R\).

Sur \(\C\), on a :
\[
X^2+1=(X-i)(X+i).
\]
Les valeurs propres sont :
\[
i,\quad -i.
\]
Elles sont distinctes.  
Donc \(A\) est diagonalisable sur \(\C\), et donc aussi trigonalisable sur \(\C\).

\subsection*{Correction de l'exercice 5}

On a :
\[
A=
\begin{pmatrix}
1&1\\
1&1
\end{pmatrix}.
\]
Le polynôme caractéristique est :
\[
\chi_A(X)=
\det
\begin{pmatrix}
X-1&-1\\
-1&X-1
\end{pmatrix}
=
(X-1)^2-1
=
X^2-2X
=
X(X-2).
\]
Les valeurs propres sont :
\[
0
\quad\text{et}\quad
2.
\]
Elles sont distinctes, donc \(A\) est diagonalisable.

Pour \(\lambda=0\),
\[
E_0=\Vect
\begin{pmatrix}
1\\-1
\end{pmatrix}.
\]
Pour \(\lambda=2\),
\[
E_2=\Vect
\begin{pmatrix}
1\\1
\end{pmatrix}.
\]
On pose :
\[
P=
\begin{pmatrix}
1&1\\
-1&1
\end{pmatrix},
\qquad
D=
\begin{pmatrix}
0&0\\
0&2
\end{pmatrix}.
\]
Alors :
\[
A=PDP^{-1}.
\]
Comme :
\[
P^{-1}=
\frac12
\begin{pmatrix}
1&-1\\
1&1
\end{pmatrix},
\]
on obtient, pour \(n\geq1\),
\[
A^n=PD^nP^{-1}
=
2^{n-1}
\begin{pmatrix}
1&1\\
1&1
\end{pmatrix}.
\]

Si \(X_0=(a,b)^T\), alors :
\[
A^nX_0
=
2^{n-1}
\begin{pmatrix}
a+b\\
a+b
\end{pmatrix}.
\]
Si \(a+b=0\), alors :
\[
A^nX_0=0
\]
pour tout \(n\geq1\).  
Sinon, la norme de \(A^nX_0\) croît comme \(2^{n-1}\).

%=========================================================
\newpage
\section{Fiche de synthèse finale}

\subsection*{1. Diagonalisation}

Un endomorphisme \(u\) est diagonalisable si \(E\) possède une base de vecteurs propres de \(u\).

Une matrice \(A\) est diagonalisable si :
\[
A=PDP^{-1},
\]
avec \(P\) inversible et \(D\) diagonale.

Les colonnes de \(P\) sont les vecteurs propres.  
Les coefficients diagonaux de \(D\) sont les valeurs propres correspondantes.

\subsection*{2. Critères de diagonalisation}

Critère géométrique :
\[
u\text{ diagonalisable}
\Longleftrightarrow
E=\bigoplus_{\lambda\in\Sp(u)}E_\lambda(u).
\]

Critère dimensionnel :
\[
u\text{ diagonalisable}
\Longleftrightarrow
\sum_{\lambda\in\Sp(u)}\dim(E_\lambda)=\dim(E).
\]

Critère par multiplicités :
\[
u\text{ diagonalisable}
\Longleftrightarrow
\chi_u\text{ est scindé et }\forall\lambda,\ \dim(E_\lambda)=m_\lambda.
\]

Critère rapide :
\[
\dim(E)=n,\quad u\text{ possède }n\text{ valeurs propres distinctes}
\Longrightarrow
u\text{ diagonalisable}.
\]

\subsection*{3. Polynômes annulateurs}

Si :
\[
P(u)=0,
\]
alors :
\[
\Sp(u)\subset\{\lambda:\ P(\lambda)=0\}.
\]

Si \(P\) est scindé à racines simples et annule \(u\), alors \(u\) est diagonalisable.

\subsection*{4. Applications de la diagonalisation}

Si :
\[
A=PDP^{-1},
\]
alors :
\[
A^n=PD^nP^{-1}.
\]

Si :
\[
Q\in\K[X],
\]
alors :
\[
Q(A)=PQ(D)P^{-1}.
\]

Si :
\[
D=\diag(\lambda_1,\dots,\lambda_n),
\]
alors :
\[
D^n=\diag(\lambda_1^n,\dots,\lambda_n^n),
\]
et :
\[
Q(D)=\diag(Q(\lambda_1),\dots,Q(\lambda_n)).
\]

\subsection*{5. Trigonalisation}

Un endomorphisme \(u\) est trigonalisable s'il existe une base dans laquelle sa matrice est triangulaire supérieure.

Une matrice \(A\) est trigonalisable si :
\[
A=PTP^{-1},
\]
avec \(T\) triangulaire supérieure.

\subsection*{6. Critère de trigonalisation}

\[
u\text{ trigonalisable sur }\K
\Longleftrightarrow
\chi_u\text{ est scindé sur }\K.
\]

En particulier :
\[
\forall A\in\M_n(\C),\quad A\text{ est trigonalisable sur }\C.
\]

\subsection*{7. Drapeaux stables}

Un drapeau est une suite :
\[
\{0\}=E_0\subset E_1\subset\cdots\subset E_n=E
\]
avec :
\[
\dim(E_k)=k.
\]

\[
u\text{ trigonalisable}
\Longleftrightarrow
u\text{ admet un drapeau stable}.
\]

\subsection*{8. Trace et déterminant}

Si \(A\) est trigonalisable et si ses valeurs propres comptées avec multiplicité sont :
\[
\lambda_1,\dots,\lambda_n,
\]
alors :
\[
\tr(A)=\lambda_1+\cdots+\lambda_n,
\]
et :
\[
\det(A)=\lambda_1\cdots\lambda_n.
\]

\subsection*{9. Erreurs classiques}

\begin{itemize}
    \item Confondre trigonalisable et diagonalisable.
    \item Croire que \(\chi_A\) scindé suffit pour diagonaliser.
    \item Oublier de comparer \(\dim(E_\lambda)\) et \(m_\lambda\).
    \item Mettre les valeurs propres dans \(D\) dans un ordre différent de celui des colonnes de \(P\).
    \item Oublier que le corps de base change tout : une matrice peut être diagonalisable sur \(\C\) mais pas sur \(\R\).
    \item Croire qu'une matrice triangulaire est toujours diagonalisable.
    \item Confondre polynôme annulateur scindé et polynôme annulateur scindé à racines simples.
\end{itemize}

%=========================================================
\section*{Conclusion pédagogique}
\addcontentsline{toc}{section}{Conclusion pédagogique}

La diagonalisation et la trigonalisation sont deux formes majeures de réduction.

La diagonalisation est la forme idéale :
\[
A=PDP^{-1}.
\]
Elle permet de découper l'espace en directions propres indépendantes et simplifie radicalement les calculs de puissances, de polynômes de matrices et de suites récurrentes vectorielles.

La trigonalisation est plus générale :
\[
A=PTP^{-1},
\]
avec \(T\) triangulaire supérieure.  
Elle permet de conserver les informations spectrales essentielles, en particulier les valeurs propres, la trace et le déterminant, même lorsque la diagonalisation est impossible.

Le cœur du chapitre est donc :

\[
\boxed{
\text{diagonaliser}=\text{trouver assez de vecteurs propres}
}
\]

et

\[
\boxed{
\text{trigonaliser}=\text{avoir un polynôme caractéristique scindé}
}
\]

Ces deux notions constituent la base de toute la théorie de la réduction des endomorphismes.

\end{document}